\documentclass[11pt,a4paper,fontset=fandol]{ctexart}
\usepackage[margin=18mm]{geometry}
\usepackage[hidelinks,unicode]{hyperref}
\pagestyle{empty}
\setlength{\parindent}{0pt}
\setlength{\parskip}{3pt}
\setlength{\emergencystretch}{2em}
\begin{document}
{\LARGE\bfseries 林予川}\par
产品实习\par
\href{mailto:candidate@example.com}{c\allowbreak{}a\allowbreak{}n\allowbreak{}d\allowbreak{}i\allowbreak{}d\allowbreak{}a\allowbreak{}t\allowbreak{}e\allowbreak{}@\allowbreak{}e\allowbreak{}x\allowbreak{}a\allowbreak{}m\allowbreak{}p\allowbreak{}l\allowbreak{}e\allowbreak{}.\allowbreak{}c\allowbreak{}o\allowbreak{}m}\par
\par\medskip\noindent{\large\bfseries 教育经历}\par\nobreak
\begin{minipage}{\linewidth}
\textbf{青禾学院} | 信息管理\par
{\small\raggedleft 2023.09—2027.06（预计）\par}
\end{minipage}\par
\par\medskip\noindent{\large\bfseries 项目经历}\par\nobreak
\begin{minipage}{\linewidth}
\textbf{团队文档检索原型} | 产品方案设计\par
{\small\raggedleft 2026.07—2026.09\par}
\noindent\hangindent=1em\hangafter=1\textbullet\ 主导团队文档检索原型的产品方案设计，在回答来源展示与多文档比较之间确定优先级，选择先展示来源并整理验收检查清单；服务由队友实现。\par
\end{minipage}\par
\end{document}
